\section{Tilting at Windmills}

There is a constant misunderstanding that seeks to treat metaphysics as a philosophy or system of thought, like perennialism, that can be studied and debated from the outside. Quite the contrary, it is more like an empirical science, although that data observed are in consciousness. Even better, the metaphysician is more like an explorer such as Columbus or Lewis and Clark. He explores the contents of his consciousness, discovers new things, and only then begins to codify them. \textbf{Rene Guenon} writes in his Hindu Doctrines:

\begin{quotex}
The real difficulty [in understanding metaphysics] is the mental assimilation needed to arrive at this result; there are certainly many minds that are quite incapable of it, and it is easy to gauge how far this effort transcends the scope of mere works of erudition. There is only one really profitable way of studying doctrines: in order to be understood they must be studied so to speak ``from the inside" …

\end{quotex}
This is the third dimension of history that \textbf{Julius Evola} refers to; that dimension is depth which is really interiority. Hence, to understand metaphysics or tradition, one must avoid the temptation of the demon of dialectics and try to engage that system in a living way, much as an explorer searching for an unexplored part of the mind. Hermetic meditation can be helpful in this regard: by an imaging, one tries to understand the state of mind of a writer or speaker.

\paragraph{Metaphysical Pathology}


Yet, as Guenon points out, few minds are capable of it. To answer why, we need to refer to the esoteric doctrine that what Guenon observes is necessarily so. Although there are many beings with a human form, only the few are capable of activating the dimension of spirit. This doctrine cannot be exoteric, since it would lead to all sorts of misunderstandings. Rather than looking for spiritual qualities, common men will try to separate men based on contingent and accidental qualities. Yet, the implication is there in the teaching on predestination. How else can we understand it?

Is it possible, then, to determine the inner state of another man? Perhaps some few enlightened beings can read a man's soul directly, but otherwise, keeping in mind the principle that the lesser cannot judge the higher, there are clues to notice. Since the outer is the reflection of the inner, a man's posture, voice, expression, and so on, often give away a lot. For example, it can be determined how conscious he is by whether his spirit dominates the soul which dominates the body, or, as is more likely, the other way around.

However, we can start at a more basic level. Since a man's understanding of metaphysics depends on his own self-knowledge, we can get an idea of his state of self-knowledge by the opinions he professes to believe. If his self-knowledge lacks breadth or depth, that will show up in his philosophy of life.

For example, I once heard an on-line lecture by a quite intelligent man who claimed a 160 IQ. Nevertheless, he regarded his own inner life as the product of electro-chemical reactions in the brain. Now, there may arise in some the urge to debate him and claim that consciousness and the soul really exist and are irreducible to material causation. Or, we could simply take him at his word. The man may simply have an impoverished inner life, whether through lack of effort or by the very nature of his constitution.

Another claim I sometimes hear, and even once from a professor, is that guilt is a Jewish invention and has no part in the Western tradition. To debate that directly is a fool's errand. The way to deal with it is by way of a phenomenology of guilt. We observe the experience in our consciousness and see how it arises. If we see it associated with the forces of chaos or neurotic instincts or false charges whose goal is manipulation, then we transcend them. Otherwise, guilt may reveal something fundamental about us. For example, \textbf{Martin Heidegger} claims to have (re)discovered a primordial sense of guilt. Now, it may very well be the case that the professor in his own consciousness cannot detect any sense of guilt. Unfortunately, this is not indicative of any higher state of mind. Rather, it is a psychological pathology which goes by the name of \textbf{sociopathy}, so don't trust such a man with your money, your daughter, or your reputation.

A third example, this time indicating an acute form of \textbf{paranoia}, is the not infrequent claim that ``Jews" are implanting Semitic ideas in our minds. Unfortunately, there is little specificity, so these alien ideas are simply ideas that the claimant doesn't like. Now, on \textbf{Terence}'s principle that \emph{nothing human is alien to me}, we need to look within ourselves for the source of such ideas; after all, a virus cannot take hold in an organism without the appropriate receptors. But the paranoid always experiences such ideas as coming from outside him. This is a pathology where the Persona projects undesired qualities onto the Shadow\footnote{\url{https://www.gornahoor.net/?p=1815}}.

\textbf{Otto Weininger} made an effort to define exactly what are Aryan and Semitic ideas. While some of his insights are useful, his use of those terms is not at all helpful and just leads to misunderstanding. Weininger does qualify it by separating the Platonic ideal of Aryanity and the Semitic anti-ideal from their presence in specific races or ethnic groups. We prefer to refer to the forces of Order and the forces of Chaos, which we have defined on too many posts to link to. We observe them within ourselves. To the extent that the forces of chaos exert influence on us, they are a privation; that is, a force that by our nature we should overcome, but we have not yet developed the inner strength to oppose it. That is our real battle, not \emph{tilting at windmills}.



\flrightit{Posted on 2011-09-19 by Cologero }

\begin{center}* * *\end{center}

\begin{footnotesize}\begin{sffamily}



\texttt{Cologero on 2011-09-19 at 12:39 said: }

There is a world of difference between transcending guilt and not being aware of it at all.

We see a similar thing in Buddhism in the West, which attracts a lot of weak men. That is because they hear doctrines like the ``no-self", and then, because they have such a weak sense of self within, they assume they are close to enlightened Buddhahood. But it is only the man with a strong self who is capable of transcending that self who is a Buddha; not at all, the man who never had any sense of self at all.


\hfill


\end{sffamily}\end{footnotesize}
